\begin{quote}
AIC is likely one of the most rewarding courses I've attended at NTNU.
It gave me a lot of valuable knowledge on different types of circuits,
IC design workflows and open source EDA tools that I greatly appreciate.
It is also one of the most challenging courses, due to the amount of
effort and time I had to spend in order to figure things out. - Tord,
AIC2025
\end{quote}

\begin{quote}
A fantastic project that just might turn your world upside down, push
you to re-evaluate your life choices, and stare briefly into the
existential void\ldots{} all while being deeply enjoyable and engaging!
- Domen, AIC2025
\end{quote}

\textbf{Keywords:} Temperature Sensor, Leakage, Regulator,
Specification, FOM, Milestones, Tapeout

The project will walk you through the full analog/digital design
process. From specification all the way to a finished layout, and a
potential tapeout.

The project is not easy, it's rather hard. You'll experience
frustration, despair, epic wins, epic losses, stress, collaboration, and
you will figure out whether you love analog design, or digital design or
neither.

I promise that the design project closely matches how we would develop a
circuit in industry.

I ask a lot of you on the project, as such, it accounts for 45 \% of the
grade, and is maybe the thing that you'll learn the most from.

In this document I'll go through the problem (what we're trying to
solve), and the milestones that we'll use along the way.

\section{The challenge}\label{the-challenge}\index{the challenge@The challenge}

The assignment is to Design a temperature sensor

But why? I'll try to explain.

\subsection{Systems-on-chip have complex regulator
systems}\label{systems-on-chip-have-complex-regulator-systems}

See the example in Figure 2 from Nordic Semiconductor's nRF54L15 product
specification.

VDD is the supply from the battery (1.7 V - 3.6 V). While the DECD, DECA
and DECRF are the low voltage supplies for the digital, analog and
radio.

The VREGMAIN has both a DC/DC, and a LDO. We'll learn about those in the
course. For now it's sufficient to know that the DC/DC converts power
drawn on the low supply (DECA, DECD, DECRF) to power drawn from the high
supply (VDD), while the LDO has the same current on low supply as the
high supply, but the voltage is different.

You will learn in the course that the typical systems inside VREGMAIN
are complicated, and sometimes complex, analog circuits.

From the data-sheet you'll see that the lowest power state is about 700
nA, while the highest power state is about 10 mA. The high power state
is 14 thousand times higher than the low power state!


{
\centering
\aicfig{media/nrf54L15_power.pdf}
}

\emph{Figure 2: Power system of nRF54L15 }

Those numbers are the total current consumption. That includes switching
currents from digital, analog bias currents, and leakage currents. In
modern technologies, because of the low threshold voltage, the leakage
currents can be a large part of the total current budget.

\subsection{Leakage current varies orders of magnitude over
temperature}\label{leakage-current-varies-orders-of-magnitude-over-temperature}

The sub-threshold leakage current in a MOSFET is

\[I_{leak} = I_0 e^{-V_{th}/n V_T} \left(1 - e^{-V_{ds}/V_T}\right)\]

The change in leakage as a function of temperature is rather
complicated. The \(V_T = \frac{k T}{q}\) factor is easy, but both
\(I_0\) and \(V_{th}\) have a complicated relation to temperature.

In Figure 3 you can see the leakage simulation (from
\url{http://analogicus.com/lelo_aic_sky130a/})


{
\centering
\aicfig{media/TB_LEAK.pdf}
}

\emph{Figure 3: Leakage simulation }

Based on the previous curves we could run a thought experiment.

\begin{itemize}
\item
  Assume 1 pA at 25 C, and 1 nA at 125 C, per logic cell
\item
  Assume 100 million logic cells
\item
  Leakage at 25 C =\textgreater{} 100 uA
\item
  Leakage at 125 C =\textgreater{} 100 mA !!!
\end{itemize}

\subsection{Why we would like to know the temperature on
die}\label{why-we-would-like-to-know-the-temperature-on-die}

Expanding on the thought experiment.

\begin{itemize}
\item
  Assume we use 1 \% of the load current for the regulator
\item
  At 25 C =\textgreater{} 1 uA for LDO
\item
  At 125 C =\textgreater{} 1 mA for LDO
\end{itemize}

It's insanely difficult to design a regulator that is efficient across
the full range of leakage currents at any temperature.

It would be good if we could know temperature.

\subsection{How to measure
temperature?}\label{how-to-measure-temperature}

There are a multitude of ways to make a temperature sensor. In
\citeproc{ref-tang20}{{[}1{]}} they used a leakage based digital ring
oscillator, in \citeproc{ref-jeong2014}{{[}2{]}} they used a
two-transistor MOSFET sensing element, in
\citeproc{ref-pertijs2005}{{[}3{]}} they had a more complicated
sigma-delta ADC sensing bipolar transistors.

The design of a temperature sensor is more difficult than you think. As
such, I would suggest that you don't go too crazy in your choice of
sensor. So far, none has gotten close to the finish line with a
sigma-delta ADC based sensor.

In the previous years of Advanced Integrated Circuits most groups have
chosen an architecture similar to \citeproc{ref-park2022}{{[}4{]}} Fig.
2. I would recommend you do the same, and that's what I'll target in the
milestones.

The principle of the temperature sensor is: 1) Create a current that is
proportional to temperature (Lecture 3), 2) Convert current to frequency
with a relaxation oscillator (Lecture 9). 3) Check the frequency to read
the temperature.

In Figure 4 below you can see an illustration of the temperature sensor.

A bandgap circuit is used to make a current that is proportional to
absolute temperature (\(I_{PTAT}\)) and a voltage that is complementary
to absolute temperature (\(V_{CTAT}\)). A relaxation oscillator converts
the current and voltage into a frequency (\(f_{OSC}\)). A digital
finite-state-machine and a counter converts the frequency to a digital
value that is proportional to temperature.

I've made an example temperature sensor at
\href{https://analogicus.com/lelo_temp_sky130a/}{lelo\_temp\_sky130a}.
Feel free to steal ideas, and circuits, from that design.


{
\centering
\aicfig{media/aic2026_project_analog.pdf}
}

\emph{Figure 4: Illustration of the temperature sensor }

\section{The Project}\label{the-project}\index{the project@The project}

I'm going to lead you through the design of a, to you, complicated mixed
signal circuit design. You will despair, you will not understand, but
you will learn.

In order to make the problem possible to learn, we're going to focus on
one milestone at a time. I hope that will enable you to not drown before
we get started.

An illustration of the milestones can be seen in Figure 5.

The first milestone is the design of the bandgap circuit, a pure analog
design. The second milestone is the design of the relaxation oscillator.
The third milestone is how you measure the frequency of the oscillator,
and is usually done in SystemVerilog.

The fourth milestone is optional, but you can't get an A in the course
if you don't have some points from the layout and parasitic simulation
milestone.

The fifth milestone is an individual report. I will force you to work in
groups. As such, it may be that some contribute more than others. To
ensure that the grading is fair, the report will be individual. It's OK
to share figures, tables, and so on, but the PDF shall be written by you
and you alone.

The sixth and last milestone is the tapeout on
\textless tinytapeout.com\textgreater. It's optional, and not everyone
will get to that stage.


{
\centering
\aicfig{media/aic2026_project.pdf}
}

\emph{Figure 5: Project overview }

\subsection{Specification}\label{specification}\index{specification@Specification}

The temperature sensor shall be designed to fit the specification below.

\begin{longtable}[]{@{}
  >{\raggedright\arraybackslash}p{(\linewidth - 8\tabcolsep) * \real{0.0865}}
  >{\raggedright\arraybackslash}p{(\linewidth - 8\tabcolsep) * \real{0.1731}}
  >{\raggedright\arraybackslash}p{(\linewidth - 8\tabcolsep) * \real{0.0673}}
  >{\raggedright\arraybackslash}p{(\linewidth - 8\tabcolsep) * \real{0.0577}}
  >{\raggedright\arraybackslash}p{(\linewidth - 8\tabcolsep) * \real{0.6154}}@{}}
\toprule\noalign{}
\begin{minipage}[b]{\linewidth}\raggedright
Key
\end{minipage} & \begin{minipage}[b]{\linewidth}\raggedright
Parameter
\end{minipage} & \begin{minipage}[b]{\linewidth}\raggedright
Value
\end{minipage} & \begin{minipage}[b]{\linewidth}\raggedright
Unit
\end{minipage} & \begin{minipage}[b]{\linewidth}\raggedright
Description
\end{minipage} \\
\midrule\noalign{}
\endhead
\bottomrule\noalign{}
\endlastfoot
Area & Area & \textless{} 15000 & um\^{}2 & Must fit in 161 um x 111 um
tiny tapeout 1x1 block, which is 17871 um\^{}2 \\
Tc & Conversion time & \textless{} 30 & us & Analog should only be
active for one 32768 Hz period \\
Ts & Sample period & 100 & ms & One conversion every 100 ms, so 10
samples per second \\
Ileak & Leakage current & \textless{} 1 & nA & Typical temperature (25
C) \\
Iact & Active current & \textless{} 100 & uA & Typical temperature (25
C) \\
Iavg & Average current & \textless{} 50 & nA & Active current x
conversion time/sample rate + leakage current. Typical temperature (25
C) \\
Kerrone & Accuracy 0 - 70C & +-10 & C & One temperature (25C)
calibration \\
Kerrtwo & Accuracy 0 - 70C & +-5 & C & Two temperature (25C, 85C)
calibration \\
\end{longtable}

\subsection{Figure of Merit}\label{figure-of-merit}\index{figure of merit@Figure of Merit}

A figure of merit allows us to compare circuits with different
performance specifications, and try to answer the question. Which one is
better?

The figure of merit for our temperature sensor will be

\[FOM = \left(\frac{T_{c}}{T_s}I_{act} + I_{leak}\right) K_{errtwo} \text{ [AK]}\]

The bracket is exactly \(I_{avg}\) from the table above, so the figure
of merit is \emph{average current times error} --- what the sensor costs
you multiplied by how wrong it is. Lower is better, and there is no way
to win by being cheap and inaccurate or by being accurate and expensive.

It is worth checking that the specification is self consistent before
designing to it, because specifications often are not. With the numbers
in the table,
\(\frac{30\ \mu s}{100\ ms}\times 100\ \mu A + 1\ nA = 31\) nA,
comfortably inside the 50 nA the table allows for \(I_{avg}\). Notice
also where the current goes: the analog block draws 100 uA but only for
30 us in every 100 ms, so duty cycling turns it into 30 nA and the
leakage, at 1 nA, is almost an afterthought. That ratio is why the
conversion time is specified at all.

\subsection{Grade}\label{grade}\index{grade@Grade}

Milestones are important, and Milestone 0 - 5 will count towards your
final grade!

That means, you have to start working right away.

The points have been designed such that it's impossible to get an A
without getting some points on the layout

\begin{longtable}[]{@{}
  >{\raggedright\arraybackslash}p{(\linewidth - 10\tabcolsep) * \real{0.0309}}
  >{\raggedright\arraybackslash}p{(\linewidth - 10\tabcolsep) * \real{0.0619}}
  >{\raggedright\arraybackslash}p{(\linewidth - 10\tabcolsep) * \real{0.0567}}
  >{\raggedright\arraybackslash}p{(\linewidth - 10\tabcolsep) * \real{0.3918}}
  >{\raggedright\arraybackslash}p{(\linewidth - 10\tabcolsep) * \real{0.3711}}
  >{\raggedright\arraybackslash}p{(\linewidth - 10\tabcolsep) * \real{0.0876}}@{}}
\toprule\noalign{}
\begin{minipage}[b]{\linewidth}\raggedright
Week
\end{minipage} & \begin{minipage}[b]{\linewidth}\raggedright
Deadline
\end{minipage} & \begin{minipage}[b]{\linewidth}\raggedright
Milestone
\end{minipage} & \begin{minipage}[b]{\linewidth}\raggedright
Task
\end{minipage} & \begin{minipage}[b]{\linewidth}\raggedright
Condition for more than 0 points
\end{minipage} & \begin{minipage}[b]{\linewidth}\raggedright
Possible Points
\end{minipage} \\
\midrule\noalign{}
\endhead
\bottomrule\noalign{}
\endlastfoot
4 & 2026-01-23 & M0 & Complete tutorial & Link on blackboard & 5 \\
7 & 2026-02-13 & M1 & Design a circuit that can convert a temperature
into a current and voltage & Description of the circuit on github docs &
5 \\
10 & 2026-03-06 & M2 & Design a circuit that can convert a temperature
into a frequency & Description of the circuit on github docs.
Demonstrate that it works & 10 \\
13 & 2026-03-27 & M3 & A verilog testbench that can convert a frequency
into a digital value & Description of the testbench on github docs.
Demonstrate that it works & 10 \\
16 & 2026-04-17 & M4 & Layout of your circuit & DRC/LVS/GDS passing on
github & 20 \\
18 & 2026-05-01 & M5 & Individual report & Uploaded to Inspera & 48 \\
? & & M6 & Tapeout & None & 0 \\
& & Coolness & Extra points that I may choose to award & & 10 \\
& & Total & & & 108 \\
\end{longtable}

\subsection{Milestone 0: The tutorial}\label{milestone-0-the-tutorial}\index{milestone 0: the tutorial@Milestone 0: The tutorial}

\textbf{Goal}: force you to install the tools, and get you started.

Follow \href{https://analogicus.com/aic2026/sky130nm_tutorial}{Sky130nm
Tutorial}

\textbf{Delivery}: Submit link to your github repository on blackboard

For example, my repository:
\href{http://analogicus.com/lelo_ex_sky130a/}{LELO\_EX\_SKY130A}

\textbf{The exercise will teach you the skills you need to do the
project}

\subsection{Milestone 1: The bandgap}\label{milestone-1-the-bandgap}\index{milestone 1: the bandgap@Milestone 1: The bandgap}

\textbf{Goal}: Create a circuit that can transform a temperature on the
integrated circuit to a current proportional to temperature (PTAT), and
a voltage complementary to temperature (CTAT).

For this purpose it's common to use ``Bandgap'' circuits. We'll learn
about them in the course, but if you don't want to wait then you should
read \url{https://analogicus.com/aic2026/references_and_bias} and ask me
questions in reference and bias lecture.

In the git repository for your group you'll create schematics for the
bandgap circuits, and you'll make test-benches to check that the bandgap
circuit works.

If you don't know what you should simulate and verify, it's good to have
a chat with ChatGPT. See
\url{https://chatgpt.com/share/69481b11-8830-8007-9986-c9e41d735cfc}.

Or check my test-benches at
\url{https://github.com/wulffern/lelo_temp_sky130a/tree/main/sim/LELOTEMP_BIAS_IBP}

\textbf{Delivery}: Link to your github repository with a description of
how the bandgap works.


{
\centering
\aicfig{media/l03_ptat_tikz.pdf}
}

\emph{Figure 6: PTAT current generator, where an op amp forces equal
voltages across the two bipolar branches so that \(\Delta V_{BE}\) over
\(R_1\) sets \(I_{PTAT}\)}

\subsection{Milestone 2: The
oscillator}\label{milestone-2-the-oscillator}

\textbf{Goal}: Use the PTAT current, and the CTAT voltage an create a
oscillator.

One way is to charge a capacitor with the current, and have a comparator
trigger when the voltage on the capacitor reaches a voltage (CTAT
voltage for example). When the comparator triggers, then we can reset
the capacitor.

This is similar to what group 7 did last year (one of the groups got all
the way to tapeout). One difference, though, is that group 7 used VDD/2
as the reference. I would recommend you use the CTAT voltage from the
bandgap instead. That way, the oscillation frequency is independent (to
first order) from the VDD.


{
\centering
\aicfig{media/rcosc_tikz.pdf}
}

\emph{Figure 7: Relaxation oscillator, where a current charges capacitor
\(C\) until the comparator trips against the reference voltage \(V_1\)
across \(R\), and the flip-flop resets the capacitor}

\textbf{Delivery} Link to your github repository with description on how
your oscillator works. There should be proof on how it works.

\subsection{Milestone 3: The
measurement}\label{milestone-3-the-measurement}

\textbf{Goal}: Measure the frequency of the oscillator.

In the system we can assume we have an accurate 32768 Hz clock source.
One way to find the frequency is to run the oscillator for a fixed
number of clock cycles on the 32768 Hz clock, and have a counter that
can count the output pulses.

Assume we counted 128 clock cycles over 2 clock periods of the 32768 Hz
clock. That would mean the frequency of the oscillator was approximately
2.09 MHz. Once we have the frequency we can calculate the temperature.

I would recommend that you write in verilog the system to start the
oscillator, count for a number of 32768 Hz clock cycles, and transform
the frequency into a temperature.


{
\centering
\aicfig{media/b6546506e714da05dd00b32acfe99950fcff6a2aea264115ce7296c2a71638f3.pdf}
}

\emph{Figure 8: State machine that measures the oscillator frequency,
cycling IDLE, PWRUP (count oscillator pulses), PWRDWN and CAPTURE (store
counter value)}

\textbf{Delivery}: Link to your github repository where you describe how
you measure the frequency of the oscillator.

\subsection{Milestone 4 (Optional): The physical
design}\label{milestone-4-optional-the-physical-design}

\textbf{Goal}: Do the physical layout of your oscillator, and prove that
it still works with the layout parasitics.

If you do the layout, then your design must fit within a digital 1x1
tinytapeout block. I've made a template at
\url{https://github.com/wulffern/lelo_temp_sky130a/blob/main/design/LELO_TEMP_SKY130A/tt_block_1x1_pg.mag}
that you can use. If the design does not fit within that space, then you
won't be able to tapeout.

When your design is complete, then the DRC, LVS, GDS actions should be
passing on github.


{
\centering
\aicfig{media/l00_layout.pdf}
}

\emph{Figure 9: Physical layout of a circuit in the Magic layout editor,
with the DRC error count shown as zero in the toolbar}

\textbf{Delivery}: Link to your github repository with passing GDS, DRC,
LVS actions.

\subsection{Milestone 5: The Report}\label{milestone-5-the-report}\index{milestone 5: the report@Milestone 5: the report}

\textbf{Goal}: Write a report

\textbf{Delivery}: A PDF copy of the report in Inspera. You'll all write
an individual report. The report shall be in the IEEE template.

See further details in
\url{https://analogicus.com/aic2026/how_to_write_a_project_report}

\subsection{Milestone 6 (Optional): The
Tapeout}\label{milestone-6-optional-the-tapeout}

Target TTSKY26b tapeout (June 2026) on
\url{https://tinytapeout.com/chips/}

Those students that follow the course at NTNU will be able to tapeout if
the design is complete. I've gotten \href{https://nordicsemi.com}{Nordic
Semiconductor} to sponsor the tapeout for 2026.

\section{Summary}\label{summary}

The one-page version of this chapter:

\begin{itemize}
\tightlist
\item
  The project is the course: design a block in an open PDK, simulate it
  over corners, and document the evidence
\item
  Start from the template repositories and the cic flow on day one -
  infrastructure debt compounds
\item
  Work like an engineer: git for everything, scripts over clicks, claims
  backed by simulation
\item
  The report is the deliverable; the next chapter's writing guide is not
  optional reading
\end{itemize}

\section{Would you like to know
more?}\label{would-you-like-to-know-more}

Precision CMOS temperature sensors, the book chapter to read before
starting \citeproc{ref-pertijs2005}{{[}3{]}}

Recent low-power sensors, for what the state of the art costs in area
and current \citeproc{ref-jeong2014}{{[}2{]}}
\citeproc{ref-tang20}{{[}1{]}} \citeproc{ref-park2022}{{[}4{]}}

\phantomsection\label{refs}
\begin{CSLReferences}{0}{0}
\bibitem[\citeproctext]{ref-tang20}
\CSLLeftMargin{{[}1{]} }%
\CSLRightInline{Z. Tang, Y. Fang, Z. Shi, X.-P. Yu, N. N. Tan, and W.
Pan, {``A 1770- \(\mu\) m2 leakage-based digital temperature sensor with
supply sensitivity suppression in 55-nm CMOS,''} \emph{IEEE Journal of
Solid-State Circuits}, vol. 55, no. 3, pp. 781--793, 2020, doi:
\href{https://doi.org/10.1109/JSSC.2019.2952855}{10.1109/JSSC.2019.2952855}.
}

\bibitem[\citeproctext]{ref-jeong2014}
\CSLLeftMargin{{[}2{]} }%
\CSLRightInline{S. Jeong, Z. Foo, Y. Lee, J.-Y. Sim, D. Blaauw, and D.
Sylvester, {``A fully-integrated 71 nW CMOS temperature sensor for low
power wireless sensor nodes,''} \emph{IEEE Journal of Solid-State
Circuits}, vol. 49, no. 8, pp. 1682--1693, 2014, doi:
\href{https://doi.org/10.1109/JSSC.2014.2325574}{10.1109/JSSC.2014.2325574}.
}

\bibitem[\citeproctext]{ref-pertijs2005}
\CSLLeftMargin{{[}3{]} }%
\CSLRightInline{M. A. P. Pertijs, A. Niederkorn, X. Ma, B. McKillop, A.
Bakker, and J. H. Huijsing, {``A CMOS smart temperature sensor with a
3/spl sigma/ inaccuracy of /spl plusmn/0.5/spl deg/c from -50/spl deg/c
to 120/spl deg/c,''} \emph{IEEE Journal of Solid-State Circuits}, vol.
40, no. 2, pp. 454--461, 2005, doi:
\href{https://doi.org/10.1109/JSSC.2004.841013}{10.1109/JSSC.2004.841013}.
}

\bibitem[\citeproctext]{ref-park2022}
\CSLLeftMargin{{[}4{]} }%
\CSLRightInline{J.-H. Park, J.-H. Hwang, C. Shin, and S.-J. Kim, {``A
BJT-based temperature-to-frequency converter with +- 1°c (3\(\sigma\))
inaccuracy from - 40 c to 140 c for on-chip thermal monitoring,''}
\emph{IEEE Journal of Solid-State Circuits}, vol. 57, no. 10, pp.
2909--2918, 2022, doi:
\href{https://doi.org/10.1109/JSSC.2022.3182708}{10.1109/JSSC.2022.3182708}.
}

\end{CSLReferences}
